\begin{pastexam}[label={px:ncku-108-ee-ds-10}]{成大電機所 108 學年度 問答題第 9 題}
（問答）

Given two sparse matrices as follows:
\[
\mathrm{M}_1=\begin{bmatrix}0&10&0&12\\0&0&0&0\\0&0&5&0\\15&12&0&0\end{bmatrix},\quad
\mathrm{M}_2=\begin{bmatrix}0&0&8&0\\0&0&0&23\\0&0&10&0\\20&25&0&0\end{bmatrix}
\]
\begin{enumerate}[label=(\Alph*)]
\item Please use the linked list to represent the above two sparse matrices (5 points)
\item Please compute the multiplication of two matrices, M$_1$ X M$_2$, using linked list
representations. Please explain your method in details. (10 points)
\end{enumerate}
\end{pastexam}

\begin{pastsoln}
\textbf{(A)} 用正交串列（十字串列，資結書第~5.3~節）：每個非零元素一個節點，存列號、行號、值，以及兩個指標
$\mathit{right}$（同一列的下一個非零元素）和 $\mathit{down}$（同一行的下一個非零元素）。 每一列、每一行各有
一個頭節點，同一列的節點依行號由小到大串起來，同一行的節點依列號由小到大串起來，最後一個節點指回頭節點
（環狀）。 下圖節點寫「列,行　值」，實線箭頭是 $\mathit{right}$，虛線箭頭是 $\mathit{down}$；指回頭節點的指標
和沒有元素的列（$\mathrm{M}_1$ 的第 2 列）指向自己的指標沒有畫出。

\begin{center}
\begin{tikzpicture}[baseline=(current bounding box.center),
    hd/.style={draw, fill=black!8, minimum width=0.75cm, minimum height=0.55cm, font=\small},
    en/.style={draw, thick, minimum width=1.2cm, minimum height=0.55cm, font=\scriptsize},
    rl/.style={-{Stealth[length=1.6mm]}, thick},
    dl/.style={-{Stealth[length=1.6mm]}, thick, densely dashed}]
\node[font=\small] at (0,0) {$\mathrm{M}_1$};
\node[hd] (C1) at (1.75,0) {C1};
\node[hd] (C2) at (3.50,0) {C2};
\node[hd] (C3) at (5.25,0) {C3};
\node[hd] (C4) at (7.00,0) {C4};
\node[hd] (R1) at (0,-0.90) {R1};
\node[hd] (R2) at (0,-1.80) {R2};
\node[hd] (R3) at (0,-2.70) {R3};
\node[hd] (R4) at (0,-3.60) {R4};
\node[en] (e12) at (3.50,-0.90) {1,2 \quad 10};
\node[en] (e14) at (7.00,-0.90) {1,4 \quad 12};
\node[en] (e33) at (5.25,-2.70) {3,3 \quad 5};
\node[en] (e41) at (1.75,-3.60) {4,1 \quad 15};
\node[en] (e42) at (3.50,-3.60) {4,2 \quad 12};
\draw[rl] (R1) -- (e12);
\draw[rl] (e12) -- (e14);
\draw[rl] (R3) -- (e33);
\draw[rl] (R4) -- (e41);
\draw[rl] (e41) -- (e42);
\draw[dl] (C1) -- (e41);
\draw[dl] (C2) -- (e12);
\draw[dl] (e12) -- (e42);
\draw[dl] (C3) -- (e33);
\draw[dl] (C4) -- (e14);
\end{tikzpicture}
\end{center}
\begin{center}
\begin{tikzpicture}[baseline=(current bounding box.center),
    hd/.style={draw, fill=black!8, minimum width=0.75cm, minimum height=0.55cm, font=\small},
    en/.style={draw, thick, minimum width=1.2cm, minimum height=0.55cm, font=\scriptsize},
    rl/.style={-{Stealth[length=1.6mm]}, thick},
    dl/.style={-{Stealth[length=1.6mm]}, thick, densely dashed}]
\node[font=\small] at (0,0) {$\mathrm{M}_2$};
\node[hd] (C1) at (1.75,0) {C1};
\node[hd] (C2) at (3.50,0) {C2};
\node[hd] (C3) at (5.25,0) {C3};
\node[hd] (C4) at (7.00,0) {C4};
\node[hd] (R1) at (0,-0.90) {R1};
\node[hd] (R2) at (0,-1.80) {R2};
\node[hd] (R3) at (0,-2.70) {R3};
\node[hd] (R4) at (0,-3.60) {R4};
\node[en] (e13) at (5.25,-0.90) {1,3 \quad 8};
\node[en] (e24) at (7.00,-1.80) {2,4 \quad 23};
\node[en] (e33) at (5.25,-2.70) {3,3 \quad 10};
\node[en] (e41) at (1.75,-3.60) {4,1 \quad 20};
\node[en] (e42) at (3.50,-3.60) {4,2 \quad 25};
\draw[rl] (R1) -- (e13);
\draw[rl] (R2) -- (e24);
\draw[rl] (R3) -- (e33);
\draw[rl] (R4) -- (e41);
\draw[rl] (e41) -- (e42);
\draw[dl] (C1) -- (e41);
\draw[dl] (C2) -- (e42);
\draw[dl] (C3) -- (e13);
\draw[dl] (e13) -- (e33);
\draw[dl] (C4) -- (e24);
\end{tikzpicture}
\end{center}

$\mathrm{M}_1$ 有 $5$ 個非零元素：$(1,2)\,10$、$(1,4)\,12$、$(3,3)\,5$、$(4,1)\,15$、$(4,2)\,12$；$\mathrm{M}_2$ 有 $5$ 個：
$(1,3)\,8$、$(2,4)\,23$、$(3,3)\,10$、$(4,1)\,20$、$(4,2)\,25$。 每個矩陣是 $5$ 個元素節點加 $8$ 個頭節點，不必存
$11$ 個 $0$。 （只用「每列一條串列」也可以表示，但 (B) 要沿著 $\mathrm{M}_2$ 的行走，有 $\mathit{down}$ 指標
比較方便。）

\textbf{(B)} $\mathrm{P} = \mathrm{M}_1 \times \mathrm{M}_2$，$\mathrm{P}[i][j] = \sum_k \mathrm{M}_1[i][k]\,\mathrm{M}_2[k][j]$。
只有 $\mathrm{M}_1[i][k]$ 和 $\mathrm{M}_2[k][j]$ \textbf{都不是 $0$} 的 $k$ 才有貢獻。 $\mathrm{M}_1$ 第 $i$ 列的串列依
行號 $k$ 排好，$\mathrm{M}_2$ 第 $j$ 行的串列依列號 $k$ 排好，所以算 $\mathrm{P}[i][j]$ 就是把這兩條排好的串列
像合併一樣同時往前走，找出 $k$ 相同的節點相乘再加起來（和資結書第~5.3~節多項式相加比較指數的走法一樣）：

\begin{pseudo}{正交串列的矩陣乘法}
\Procedure{Sparse-Mult}{$M_1, M_2, n$}
    \State 建立空的正交串列 $P$（$n$ 個列頭、$n$ 個行頭）
    \For{$i \gets 1$ \To $n$}
        \For{$j \gets 1$ \To $n$}
            \State $p \gets M_1.\mathit{row}[i].\mathit{right}$；$q \gets M_2.\mathit{col}[j].\mathit{down}$；$s \gets 0$
            \While{$p$ 不是列頭 \AND $q$ 不是行頭}
                \If{$p.\mathit{col} = q.\mathit{row}$}
                    \State $s \gets s + p.\mathit{val} \cdot q.\mathit{val}$
                    \State $p \gets p.\mathit{right}$；$q \gets q.\mathit{down}$
                \ElsIf{$p.\mathit{col} < q.\mathit{row}$}
                    \State $p \gets p.\mathit{right}$
                \Else
                    \State $q \gets q.\mathit{down}$
                \EndIf
            \EndWhile
            \If{$s \ne 0$}
                \State 新節點 $(i, j, s)$ 接在 $P$ 第 $i$ 列的尾端和第 $j$ 行的尾端
            \EndIf
        \EndFor
    \EndFor
\EndProcedure
\end{pseudo}
$j$ 由小到大，所以第 $i$ 列的新節點本來就依行號排好；$i$ 由小到大，第 $j$ 行也依列號排好。 每一列、每一行
各記一個尾指標，接上去是 $O(1)$。

\textbf{計算}（只列兩條串列都不空的 $(i, j)$；$\mathrm{M}_1$ 第 2 列是空的，整列結果都是 $0$）：
\begin{center}
\small
\begin{tabular}{cllll}
\toprule
$(i, j)$ & $\mathrm{M}_1$ 第 $i$ 列 $(k, \text{值})$ & $\mathrm{M}_2$ 第 $j$ 行 $(k, \text{值})$ & 相同的 $k$ & $\mathrm{P}[i][j]$ \\
\midrule
$(1,1)$ & $(2,10), (4,12)$ & $(4,20)$ & $4$ & $12 \cdot 20 = 240$ \\
$(1,2)$ & $(2,10), (4,12)$ & $(4,25)$ & $4$ & $12 \cdot 25 = 300$ \\
$(1,3)$ & $(2,10), (4,12)$ & $(1,8), (3,10)$ & 無 & $0$ \\
$(1,4)$ & $(2,10), (4,12)$ & $(2,23)$ & $2$ & $10 \cdot 23 = 230$ \\
$(3,1)$ & $(3,5)$ & $(4,20)$ & 無 & $0$ \\
$(3,2)$ & $(3,5)$ & $(4,25)$ & 無 & $0$ \\
$(3,3)$ & $(3,5)$ & $(1,8), (3,10)$ & $3$ & $5 \cdot 10 = 50$ \\
$(3,4)$ & $(3,5)$ & $(2,23)$ & 無 & $0$ \\
$(4,1)$ & $(1,15), (2,12)$ & $(4,20)$ & 無 & $0$ \\
$(4,2)$ & $(1,15), (2,12)$ & $(4,25)$ & 無 & $0$ \\
$(4,3)$ & $(1,15), (2,12)$ & $(1,8), (3,10)$ & $1$ & $15 \cdot 8 = 120$ \\
$(4,4)$ & $(1,15), (2,12)$ & $(2,23)$ & $2$ & $12 \cdot 23 = 276$ \\
\bottomrule
\end{tabular}
\end{center}
走法舉例，$(1,1)$：$p$ 在 $(1,2)$、$q$ 在 $(4,1)$，比較 $k$：$2 < 4$，$p$ 前進到 $(1,4)$；$4 = 4$，相乘 $12 \cdot 20$，
兩個都前進，走到頭節點，結束。 $(1,3)$：$p$ 的 $k = 2$、$q$ 的 $k = 1$，$q$ 前進；$2 < 3$，$p$ 前進；$4 > 3$，$q$
前進到頭節點，結束，$s = 0$，不建節點。

結果：
\[
\mathrm{M}_1 \times \mathrm{M}_2 = \begin{bmatrix}240&300&0&230\\0&0&0&0\\0&0&50&0\\0&0&120&276\end{bmatrix},
\]
用正交串列存是 $6$ 個元素節點：$(1,1)\,240$、$(1,2)\,300$、$(1,4)\,230$ 串在第 1 列，$(3,3)\,50$ 在第 3 列，
$(4,3)\,120$、$(4,4)\,276$ 在第 4 列；行的方向，第 3 行是 $(3,3) \to (4,3)$，第 4 行是 $(1,4) \to (4,4)$。

\textbf{複雜度}。 $n \times n$ 的矩陣，$\mathrm{M}_1$ 第 $i$ 列有 $r_i$ 個非零、$\mathrm{M}_2$ 第 $j$ 行有 $c_j$ 個非零時，
算 $\mathrm{P}[i][j]$ 的合併最多走 $r_i + c_j$ 步，總共
$\sum_i \sum_j (1 + r_i + c_j) = n^2 + n t_1 + n t_2$，$t_1$、$t_2$ 是兩個矩陣的非零個數，也就是 $O(n^2 + n(t_1 + t_2))$。
直接用二維陣列乘是 $\Theta(n^3)$，矩陣越稀疏，省得越多。 （另一種寫法是對 $\mathrm{M}_1$ 第 $i$ 列的每個
$(i, k, a)$，走 $\mathrm{M}_2$ 第 $k$ 列，把 $a \cdot b$ 累加到一個長度 $n$ 的暫存列，最後把非零的收成第 $i$ 列；
結果相同。）
\end{pastsoln}
